\chapter{量子計算の基礎}
\label{chap:quantum-basics}

本章では，HHLアルゴリズムを説明するために必要な
量子計算の基礎を準備する\cite{shimada2020}．
量子状態は有限個の複素数を使って表す．
測定については，射影と呼ばれる線形写像を使う場合だけを扱う．
量子状態，状態の変化，複数の量子ビット，測定の順に説明する．
複数の量子レジスタを扱うために，テンソル積も定義する．
さらに，基底，内積，線形写像の合成と共役転置に関する性質を証明する．

\section{量子ビットと量子状態}
\label{sec:qubit-state}

\subsection{ブラケット記法}

複素数を成分にもつ有限次元のベクトル空間を$\mathcal H$とする．
空間$\mathcal H$には内積が定められているものとする．
$\mathcal H$のベクトルを
\begin{equation}
  \ket{\psi}
  \coloneq
  \begin{pmatrix}
    \psi_0&\psi_1&\cdots&\psi_{d-1}
  \end{pmatrix}^{\mathsf T}
  \label{eq:ket-coordinate}
\end{equation}
と書き，これを\textbf{ケット}という．
ケット$\ket{\psi}$の各成分を複素共役にして横に並べたものを
$\bra{\psi}:=\ket{\psi}^{\dagger}$と書き，\textbf{ブラ}という．
ベクトル$\ket{\psi}$と$\ket{\phi}$の内積は
$\braket{\phi|\psi}:=\bra{\phi}\ket{\psi}$と書く．座標を使えば，
\begin{equation}
  \braket{\phi|\psi}=\sum_{j=0}^{d-1}\overline{\phi_j}\psi_j
\end{equation}
である．内積から定まるノルムを
$\lVert\ket{\psi}\rVert:=\sqrt{\braket{\psi|\psi}}$と書く．
また，$\ket{\psi}\bra{\phi}$は線形写像を表す．
この写像は$\ket{v}$を$\braket{\phi|v}\ket{\psi}$へ移す．

% \begin{prop}
% 正規直交基底$\{\ket{e_j}\}_{j=0}^{d-1}$に対し，任意の$\ket{\psi}\in\mathcal H$は
% \begin{equation}
%   \ket{\psi}=\sum_{j=0}^{d-1}\braket{e_j|\psi}\ket{e_j}
%   \label{eq:orthonormal-expansion}
% \end{equation}
% と一意的に表される．
% \end{prop}
% \begin{proof}
% 基底であるから$\ket{\psi}=\sum_j c_j\ket{e_j}$と一意的に表される．
% 左から$\bra{e_k}$を作用させると
% \[
%   \braket{e_k|\psi}
%   =\sum_jc_j\braket{e_k|e_j}
%   =\sum_jc_j\delta_{kj}=c_k
% \]
% である．したがって$c_j=\braket{e_j|\psi}$となる．
% \end{proof}

\subsection{量子状態の公理}

% \begin{defi}[純粋状態]
% ほかの系から影響を受けない量子系を考える．
% その純粋状態は，状態空間$\mathcal H$の長さ1のベクトル$\ket{\psi}$で表す．
% つまり，$\braket{\psi|\psi}=1$とする．
% \end{defi}

ベクトルをその長さで割り，長さを1にそろえることを\textbf{規格化}という．

1量子ビットの状態空間は$\mathbb C^2$である．計算基底を
\begin{equation}
  \ket{0}=\begin{pmatrix}1\\0\end{pmatrix},
  \qquad
  \ket{1}=\begin{pmatrix}0\\1\end{pmatrix}
\end{equation}
と定めると，一般の1量子ビット状態は
\begin{equation}
  \ket{\psi}=\alpha\ket{0}+\beta\ket{1},
  \qquad |\alpha|^2+|\beta|^2=1
  \label{eq:single-qubit-state}
\end{equation}
と表される．係数$\alpha,\beta$を\textbf{確率振幅}という．
実数$\gamma$に対して$\ket{\psi}$と$e^{\mathrm{i}\gamma}\ket{\psi}$は
同じ物理状態を表す．この因子を\textbf{グローバル位相}という．
これに対して，$\alpha$と$\beta$の間の位相の差は，後の演算と測定結果に影響する．
この差を\textbf{相対位相}という．

\subsection{射影測定の公理}

\begin{defi}[直交射影]
線形写像$P:\mathcal H\to\mathcal H$が$P^2=P$かつ$P^\dagger=P$を満たすとき，
$P$を直交射影という．
\end{defi}

\begin{defi}[射影測定]
測定結果の集合を$M$とする．各結果$m\in M$に対して直交射影$P_m$を用意し，
\begin{equation}
  P_mP_{m'}=\delta_{mm'}P_m,
  \qquad \sum_{m\in M}P_m=I
  \label{eq:projective-measurement-completeness}
\end{equation}
を満たすとする．状態$\ket{\psi}$の測定で結果$m$を得る確率は
\begin{equation}
  p(m)=\bra{\psi}P_m\ket{\psi}
      =\lVert P_m\ket{\psi}\rVert^2
  \label{eq:born-rule-projector}
\end{equation}
である．この確率をボルン則という．確率が$p(m)>0$なら，測定後の状態は
\begin{equation}
  \ket{\psi_m}=\frac{P_m\ket{\psi}}{\sqrt{p(m)}}
  \label{eq:post-measurement-state}
\end{equation}
であるとする．
\end{defi}

\begin{prop}
射影測定の各結果の確率は非負で総和は1であり，
式\eqref{eq:post-measurement-state}の状態は規格化されている．
\end{prop}
\begin{proof}
式\eqref{eq:born-rule-projector}はノルムの二乗であるから$p(m)\geq0$である．
完全性関係と状態の規格化より
\[
  \sum_m p(m)
  =\bra{\psi}\left(\sum_mP_m\right)\ket{\psi}
  =\bra{\psi}I\ket{\psi}=1
\]
を得る．さらに$P_m^\dagger P_m=P_m$より
\[
  \braket{\psi_m|\psi_m}
  =\frac{\bra{\psi}P_m^\dagger P_m\ket{\psi}}{p(m)}=1
\]
となる．
\end{proof}

計算基底で測定する場合は，$P_0=\ket{0}\bra{0}$，$P_1=\ket{1}\bra{1}$とする．
式\eqref{eq:single-qubit-state}に対し，$p(0)=|\alpha|^2$，
$p(1)=|\beta|^2$である．

\begin{prop}
状態$\ket{\psi}$と$e^{\mathrm{i}\gamma}\ket{\psi}$は，すべての射影測定で同じ確率を与える．
\end{prop}
\begin{proof}
任意の$m$について
\[
  e^{-\mathrm{i}\gamma}\bra{\psi}P_m
  e^{\mathrm{i}\gamma}\ket{\psi}
  =\bra{\psi}P_m\ket{\psi}
\]
であるため，測定確率は大域位相に依存しない．
\end{proof}

\begin{ex}
状態$\ket{\psi}=(\ket{0}+\mathrm{i}\ket{1})/\sqrt2$を計算基底で測定すると，
結果$0$と$1$はともに確率$1/2$で得られる．
\end{ex}

\section{量子ゲートと複数量子ビット}
\label{sec:gates-composite}

\subsection{ユニタリ変換と1量子ビットゲート}

量子状態に行列を掛けると，別の量子状態へ変化する．
ただし，変化の前後でベクトルの長さが変わってはならない．
この条件を満たす行列または線形写像をユニタリという．

\begin{defi}[ユニタリ写像]
内積空間$\mathcal H$上の線形写像$U$が
$U^\dagger U=UU^\dagger=I$を満たすとき，$U$をユニタリ写像という．
\end{defi}

量子系が測定されず，外部から影響を受けない間，
その状態変化はユニタリ写像で表されるとする．
したがって，初期状態が$\ket{\psi}$なら，変化後の状態は$U\ket{\psi}$である．

\begin{prop}
ユニタリ写像$U$は内積とノルムを保存し，逆写像は$U^{-1}=U^\dagger$である．
\end{prop}
\begin{proof}
任意の$\ket{\phi},\ket{\psi}$について
\[
  \braket{U\phi|U\psi}
  =\bra{\phi}U^\dagger U\ket{\psi}
  =\braket{\phi|\psi}
\]
である．両ベクトルを$\ket{\phi}=\ket{\psi}$とすればノルムの保存が従う．
また，$U^\dagger U=UU^\dagger=I$である．
したがって，$U^\dagger$は$U$の逆写像である．
\end{proof}

HHLまでに用いる基本的な1量子ビットゲートを
\begin{align}
  X&=\begin{pmatrix}0&1\\1&0\end{pmatrix},&
  Y&=\begin{pmatrix}0&-\mathrm{i}\\\mathrm{i}&0\end{pmatrix},&
  Z&=\begin{pmatrix}1&0\\0&-1\end{pmatrix},\\
  H&=\frac1{\sqrt2}\begin{pmatrix}1&1\\1&-1\end{pmatrix}
\end{align}
で定める．直接の行列計算により
$X^\dagger X=Y^\dagger Y=Z^\dagger Z=H^\dagger H=I$であるから，
これらはすべてユニタリ写像である．特に
\begin{equation}
  H\ket{0}=\frac{\ket{0}+\ket{1}}{\sqrt2},
  \qquad
  H\ket{1}=\frac{\ket{0}-\ket{1}}{\sqrt2}
  \label{eq:hadamard-action}
\end{equation}
である．

等式$Y^2=I$を用いると，$y$軸回りの回転は
\begin{align}
  R_y(\theta)
  &:=\exp\left(-\frac{\mathrm{i}\theta}{2}Y\right)\\
  &=\cos\frac{\theta}{2}I-\mathrm{i}\sin\frac{\theta}{2}Y
   =\begin{pmatrix}
      \cos(\theta/2)&-\sin(\theta/2)\\
      \sin(\theta/2)& \cos(\theta/2)
    \end{pmatrix}
  \label{eq:ry-definition}
\end{align}
となる．第2の等号は指数関数の級数を偶数次と奇数次に分けて得られる．
式\eqref{eq:ry-definition}の実行列表示において$c=\cos(\theta/2)$，
$s=\sin(\theta/2)$とおけば，
\[
 R_y(\theta)^\dagger R_y(\theta)
 =\begin{pmatrix}c&s\\-s&c\end{pmatrix}
  \begin{pmatrix}c&-s\\s&c\end{pmatrix}
 =\begin{pmatrix}c^2+s^2&0\\0&c^2+s^2\end{pmatrix}=I
\]
であるから，$R_y(\theta)$もユニタリ写像である．
したがって
\begin{equation}
  R_y(\theta)\ket{0}
  =\cos\frac{\theta}{2}\ket{0}
   +\sin\frac{\theta}{2}\ket{1}
  \label{eq:ry-on-zero}
\end{equation}
である．

\subsection{テンソル積と複合系}

二つの量子ビットを同時に扱うには，二つのベクトル空間を組み合わせる必要がある．
この組合せに使う演算がテンソル積である．
まず，資料\cite[73--75頁]{tsukano2023}に従い，
行列のテンソル積をブロック行列として定義する．

\begin{defi}[行列のテンソル積]
$m\times n$行列$A=(a_{ij})$と$r\times s$行列$B$に対し，
\begin{equation}
 A\otimes B
 :=
 \begin{pmatrix}
  a_{11}B & a_{12}B & \cdots & a_{1n}B\\
  a_{21}B & a_{22}B & \cdots & a_{2n}B\\
  \vdots  & \vdots  & \ddots & \vdots\\
  a_{m1}B & a_{m2}B & \cdots & a_{mn}B
 \end{pmatrix}
 \label{eq:matrix-tensor-product-definition}
\end{equation}
を$A$と$B$のテンソル積という．これは$mr\times ns$行列である．
\end{defi}

列ベクトル$v=(v_1,\ldots,v_m)^{\mathsf T}$と
$w=(w_1,\ldots,w_r)^{\mathsf T}$も行列とみなせるため，
式\eqref{eq:matrix-tensor-product-definition}から
\begin{equation}
 v\otimes w
 =
 \begin{pmatrix}
  v_1w\\v_2w\\\vdots\\v_mw
 \end{pmatrix}
 \in\mathbb C^{mr}
 \label{eq:vector-tensor-product-coordinate}
\end{equation}
となる．つまり，各成分$v_i$をベクトル$w$全体に掛け，
その結果を上から順に並べる．

次に，一般の有限次元ベクトル空間で使う形を定める．

% \begin{defi}[有限次元ベクトル空間のテンソル積]
% 複素ベクトル空間$V,W$の基底を，それぞれ
% $\{e_i\}_{i=1}^{m}$，$\{f_j\}_{j=1}^{r}$とする．
% $V\otimes W$を，記号$e_i\otimes f_j$を基底にもつベクトル空間とする．
% また，$v=\sum_i a_ie_i$，$w=\sum_j b_jf_j$に対して
% \begin{equation}
%  v\otimes w
%  :=\sum_{i=1}^{m}\sum_{j=1}^{r}a_ib_j(e_i\otimes f_j)
%  \label{eq:vector-space-tensor-product-definition}
% \end{equation}
% と定める．
% \end{defi}

% この定義は，座標を選んだときの式\eqref{eq:vector-tensor-product-coordinate}と一致する．
特に，テンソル積は
\begin{align}
  (v_1+v_2)\otimes w&=v_1\otimes w+v_2\otimes w,\\
  v\otimes(w_1+w_2)&=v\otimes w_1+v\otimes w_2,\\
  (cv)\otimes w&=v\otimes(cw)=c(v\otimes w)
\end{align}
を満たす．
% $v\otimes w$の形のベクトルを純テンソルという．
% 一般のテンソルは純テンソルの有限個の和で表される．
% ただし，一つの純テンソルで表せるとは限らない．

% \begin{prop}
% ベクトルの族$\{e_i\}_{i=1}^{r}$を$V$の基底，$\{f_j\}_{j=1}^{s}$を$W$の基底とすると，
% $\{e_i\otimes f_j\}_{i,j}$は$V\otimes W$の基底である．
% したがって$\dim(V\otimes W)=rs$である．
% \end{prop}
% \begin{proof}
% 定義より，$e_i\otimes f_j$は$V\otimes W$の基底である．
% その個数は$rs$個であるため，$\dim(V\otimes W)=rs$となる．
% また，任意のベクトル$v=\sum_i a_ie_i$，$w=\sum_jb_jf_j$について
% \[
%   v\otimes w=\sum_{i,j}a_ib_j(e_i\otimes f_j)
% \]
% と表される．これは，式\eqref{eq:vector-space-tensor-product-definition}そのものである．
% \end{proof}

% 内積空間$V,W$に対し，純テンソル上で
% \begin{equation}
%   \langle v_1\otimes w_1,v_2\otimes w_2\rangle
%   :=\langle v_1,v_2\rangle\langle w_1,w_2\rangle
%   \label{eq:tensor-inner-product}
% \end{equation}
% と定め，線形性によって$V\otimes W$全体へ拡張する．
% もとの基底が正規直交基底なら
% \[
%  \langle e_i\otimes f_j,e_k\otimes f_l\rangle
%  =\delta_{ik}\delta_{jl}
% \]
% であるから，積基底も正規直交基底である．

% \begin{defi}[線形写像のテンソル積]
% 線形写像$A:V\to V'$，$B:W\to W'$に対し，
% $(A\otimes B)(v\otimes w)=Av\otimes Bw$を満たす線形写像を
% $A\otimes B$と定める．
% \end{defi}

% \begin{prop}
% 写像の定義域と値域が適合するとき
% \begin{align}
%  (A\otimes B)(C\otimes D)&=AC\otimes BD,\label{eq:tensor-product-composition}\\
%  (A\otimes B)^\dagger&=A^\dagger\otimes B^\dagger
%  \label{eq:tensor-product-adjoint}
% \end{align}
% が成り立つ．したがって，$A,B$がユニタリなら$A\otimes B$もユニタリである．
% これらは，行列のテンソル積の基本的な性質である
% \cite[75頁，定理2.8]{tsukano2023}．
% \end{prop}
% \begin{proof}
% 純テンソル$v\otimes w$に作用させると
% \[
%  (A\otimes B)(C\otimes D)(v\otimes w)
%  =ACv\otimes BDw=(AC\otimes BD)(v\otimes w)
% \]
% である．純テンソルが空間を生成するため式\eqref{eq:tensor-product-composition}が従う．
% また純テンソルについて
% \begin{align*}
%  \langle v'\otimes w',(A\otimes B)(v\otimes w)\rangle
%  &=\langle v',Av\rangle\langle w',Bw\rangle\\
%  &=\langle A^\dagger v',v\rangle\langle B^\dagger w',w\rangle\\
%  &=\langle(A^\dagger\otimes B^\dagger)(v'\otimes w'),v\otimes w\rangle
% \end{align*}
% であり，線形性から式\eqref{eq:tensor-product-adjoint}を得る．
% 写像$A,B$がユニタリなら両式より
% $(A\otimes B)^\dagger(A\otimes B)=I\otimes I$である．
% \end{proof}

複数の量子系を合わせたものを複合量子系という．
複合量子系の状態空間は，各量子系の状態空間のテンソル積で表す．
たとえば2量子ビットの計算基底は
\[
 \ket{00},\ket{01},\ket{10},\ket{11},
 \qquad \ket{ij}:=\ket{i}\otimes\ket{j}
\]
である．本論文では左側の因子を上位の量子ビットとし，
$\ket{x_{n-1}\cdots x_0}$を整数$\sum_{j=0}^{n-1}2^jx_j$に対応させる．

\subsection{複数量子ビットに作用するゲート}

\begin{defi}[制御ユニタリ]
一方の量子ビットを制御量子ビットとする．
また，状態空間$\mathcal K$で表される側を標的とする．このとき，
\begin{equation}
 C(U)=\ket{0}\bra{0}\otimes I_{\mathcal K}
      +\ket{1}\bra{1}\otimes U
 \label{eq:controlled-unitary}
\end{equation}
を制御$U$演算という．
\end{defi}

\begin{prop}
写像$U$がユニタリなら$C(U)$もユニタリである．
\end{prop}
\begin{proof}
射影を$P_0=\ket{0}\bra{0}$，$P_1=\ket{1}\bra{1}$とおくと，
$P_iP_j=\delta_{ij}P_i$，$P_0+P_1=I$である．よって
\begin{align*}
 C(U)^\dagger C(U)
 &=(P_0\otimes I+P_1\otimes U^\dagger)
   (P_0\otimes I+P_1\otimes U)\\
 &=P_0\otimes I+P_1\otimes U^\dagger U
  =(P_0+P_1)\otimes I=I
\end{align*}
である．同様に$C(U)C(U)^\dagger=I$も得られる．
\end{proof}

標的に作用する演算を$U=X$としたものがCNOTゲートである．
また，制御量子ビットが$\alpha\ket{0}+\beta\ket{1}$であり，
標的の状態が$\ket{\psi}$であるとする．線形性より
\begin{equation}
 C(U)(\alpha\ket{0}+\beta\ket{1})\ket{\psi}
 =\alpha\ket{0}\ket{\psi}+\beta\ket{1}U\ket{\psi}
 \label{eq:controlled-unitary-superposition}
\end{equation}
となる．この式は位相推定と制御回転で繰り返し用いる．

\section{量子回路の表現}
\label{sec:circuit-notation}

量子回路は，量子状態に加える演算と測定を回路図で表したものである．
本論文では量子線を横向きに描き，時間は左から右へ進むものとする．
1量子ビットゲートは量子線上の箱，制御ユニタリは制御点と標的の箱を結ぶ線，
測定はメータ形の記号で表す．複数量子ビットを一つのまとまりとして扱うとき，
この量子ビットのまとまりをレジスタという．

回路上で$U_1,U_2,\ldots,U_k$が左からこの順に現れる場合，
入力$\ket{\psi}$に対する出力は
\begin{equation}
 U_k\cdots U_2U_1\ket{\psi}
 \label{eq:circuit-composition-order}
\end{equation}
である．行列の積では，右端の演算から順に状態へ作用する．

\begin{prop}
ユニタリ回路$W=U_k\cdots U_1$の逆回路は
$W^\dagger=U_1^\dagger\cdots U_k^\dagger$である．
\end{prop}
\begin{proof}
行列の積の共役転置は$(AB)^\dagger=B^\dagger A^\dagger$であるから，
$(U_k\cdots U_1)^\dagger=U_1^\dagger\cdots U_k^\dagger$となる．
各$U_j$はユニタリなので$W^\dagger W=I$である．
したがって，逆回路を作るにはゲートの順を逆にし，
各ゲートをその共役転置に置き換えればよい．
\end{proof}

たとえば2量子ビットに$H\otimes I$を作用させた後，
第1量子ビットを制御とするCNOTを作用させる回路では，入力$\ket{00}$は
\[
 \ket{00}
 \xmapsto{H\otimes I}
 \frac{\ket{00}+\ket{10}}{\sqrt2}
 \xmapsto{\mathrm{CNOT}}
 \frac{\ket{00}+\ket{11}}{\sqrt2}
\]
と変化する．このように，本論文では重要な回路について，
回路図だけでなく各段階のケットも併記する．

\section{量子もつれと複合系の測定}
\label{sec:entanglement-measurement}

\begin{defi}[積状態と量子もつれ]
複合系$\mathcal H_A\otimes\mathcal H_B$の純粋状態$\ket{\Psi}$が
$\ket{u}\otimes\ket{v}$と表されるとき，$\ket{\Psi}$を積状態という．
積状態として表せない純粋状態をもつれ状態という．
\end{defi}

\begin{prop}
ベル状態$\ket{\Phi^+}=(\ket{00}+\ket{11})/\sqrt2$はもつれ状態である．
\end{prop}
\begin{proof}
積状態であると仮定し，
\[
 \ket{\Phi^+}=(a\ket0+b\ket1)(c\ket0+d\ket1)
 =ac\ket{00}+ad\ket{01}+bc\ket{10}+bd\ket{11}
\]
と書く．係数比較より$ac=bd=1/\sqrt2$かつ$ad=bc=0$である．
条件$ac\ne0$より$a,c\ne0$である．等式$ad=0$から$d=0$となるが，
これは$bd=1/\sqrt2$に反する．よって積状態ではない．
\end{proof}

複合系の第1部分だけを射影$P_m$で測定する操作は，
全体系では$P_m\otimes I$によって表される．
状態を
\begin{equation}
 \ket{\Psi}=\sum_{i,j}\alpha_{ij}\ket{i}_A\ket{j}_B
 \label{eq:general-bipartite-state}
\end{equation}
とし，第1系を計算基底で測定する．結果$i$に対応する射影は
$P_i=\ket{i}\bra{i}\otimes I$であり，
\[
 P_i\ket{\Psi}=\sum_j\alpha_{ij}\ket{i}\ket{j}
\]
となる．積基底の正規直交性から，結果$i$の確率は
\begin{equation}
 p(i)=\sum_j|\alpha_{ij}|^2
 \label{eq:partial-measurement-probability}
\end{equation}
である．確率が$p(i)>0$のとき，測定後状態は
\begin{equation}
 \frac{1}{\sqrt{p(i)}}\sum_j\alpha_{ij}\ket{i}\ket{j}
 \label{eq:partial-post-measurement-state}
\end{equation}
となる．

ベル状態では第1量子ビットの測定結果$0,1$はともに確率$1/2$である．
結果が$0$なら全体系は$\ket{00}$，結果が$1$なら$\ket{11}$となる．
したがって，第1量子ビットの測定結果を知ると，第2量子ビットの結果も確定する．

\section{補助量子ビットと逆演算}
\label{sec:ancilla-uncomputation}

量子回路では，途中の計算結果を保持するために追加の量子ビットを用いる．
これを補助量子ビット，複数まとめたものを補助レジスタという．
補助レジスタに途中の情報が残ると，目的のレジスタと結び付いたままになる．
その結果，後の計算が正しく行えないことがある．
そこで，必要な操作を終えた後に計算を逆向きに行い，
補助レジスタを初期状態へ戻す．この操作を逆演算またはアンコンピュテーションという．

系レジスタの計算基底状態$\ket{x}$に対して
\begin{equation}
 U\ket{x}\ket0=\ket{x}\ket{f(x)}
 \label{eq:ancilla-computation}
\end{equation}
となるユニタリ演算$U$を考える．
次に，$f(x)$の値に応じて変わる制御演算$V$を別のレジスタに加える．
その後に$U^\dagger$を作用させると
\begin{align}
 \sum_x\alpha_x\ket{x}\ket0\ket\eta
 &\xmapsto{U}
 \sum_x\alpha_x\ket{x}\ket{f(x)}\ket\eta\\
 &\xmapsto{V}
 \sum_x\alpha_x\ket{x}\ket{f(x)}V_{f(x)}\ket\eta\\
 &\xmapsto{U^\dagger}
 \sum_x\alpha_x\ket{x}\ket0V_{f(x)}\ket\eta
 \label{eq:uncomputation-superposition}
\end{align}
となる．最後の等号は，$U^\dagger U=I$と線形性による．
標的の状態$V_{f(x)}\ket\eta$には$f(x)$に応じた変化が残る．
一方，計算用レジスタは$\ket0$に戻る．

ただし，$U^\dagger$を作用させる前に$\ket{x}$または$\ket{f(x)}$を測定したり，
式\eqref{eq:ancilla-computation}の対応を壊す操作を加えたりすると，
上の逆演算は成立しない．HHLでは，固有値に応じた制御回転を終えた後，
逆量子位相推定によって位相レジスタを初期状態へ戻す．

\section{古典データの符号化}
\label{sec:data-encoding}

古典データを量子アルゴリズムで扱うには，データを量子状態に変える必要がある．
この変換を符号化という．
本論文では，基底符号化と振幅符号化の二つを用いる\cite{shimada2020}．

古典的な$n$ビット列$x=x_{n-1}\cdots x_0$を$n$量子ビットの計算基底状態
$\ket{x_{n-1}\cdots x_0}$へ対応させる方法を基底符号化という．
この方法は，位相推定で整数や小数をレジスタに保存するときに用いる．

\begin{defi}[振幅符号化]
非零ベクトル$\bm b=(b_0,\ldots,b_{N-1})^{\mathsf T}\in\mathbb C^N$に対し，
\begin{equation}
 \ket b=\frac1{\lVert\bm b\rVert_2}
          \sum_{j=0}^{N-1}b_j\ket j
 \label{eq:amplitude-encoding}
\end{equation}
を$\bm b$の振幅符号化という．次元$N$が2の冪でない場合は，
零成分を加えて$2^n$次元の状態空間へ埋め込む．
\end{defi}

\begin{prop}
式\eqref{eq:amplitude-encoding}で定まる$\ket b$は規格化されている．
\end{prop}
\begin{proof}
計算基底の正規直交性を用いると
\[
 \braket{b|b}
 =\frac1{\lVert\bm b\rVert_2^2}
   \sum_{j,k}\overline{b_j}b_k\braket{j|k}
 =\frac{\sum_j|b_j|^2}{\lVert\bm b\rVert_2^2}=1
\]
である．非零性の仮定$\bm b\ne\bm0$により分母は零でない．
\end{proof}

HHLへの入力は式\eqref{eq:amplitude-encoding}の$\ket b$である．
ただし，古典的なベクトル$\bm b$が与えられていても，
量子状態$\ket b$を短い時間で準備できるとは限らない．
そのため，第4章では状態を準備する方法と必要な計算量を条件として書く．

\section{制御回転による振幅の変換}
\label{sec:controlled-rotation}

値レジスタの計算基底を$\{\ket z\}$とし，
実数値関数$g$がすべての$z$について$|g(z)|\leq1$を満たすとする．
値$z$に応じて補助量子ビットを回転する演算$W_g$を
\begin{equation}
 W_g\ket z\ket0
 =\ket z\left(
   \sqrt{1-g(z)^2}\ket0+g(z)\ket1
 \right)
 \label{eq:controlled-rotation-general}
\end{equation}
と定める．

回転角を$\theta_z=2\arcsin g(z)$とおけば，
式\eqref{eq:ry-on-zero}より補助量子ビットへの$R_y(\theta_z)$によって
式\eqref{eq:controlled-rotation-general}が得られる．
値レジスタ全体では
\begin{equation}
 W_g=\sum_z\ket z\bra z\otimes R_y(2\arcsin g(z))
 \label{eq:controlled-rotation-block}
\end{equation}
と書ける．

\begin{prop}
演算$W_g$はユニタリである．また，規格化された入力
$\sum_z\alpha_z\ket z\ket0$に$W_g$を作用させて補助量子ビットを測定すると，
結果$1$を得る確率は
\begin{equation}
 p_1=\sum_z|\alpha_z|^2g(z)^2
 \label{eq:controlled-rotation-success-probability}
\end{equation}
である．成功確率が$p_1>0$のとき，値レジスタの条件付き状態は
\begin{equation}
 \frac1{\sqrt{p_1}}\sum_z\alpha_zg(z)\ket z
 \label{eq:controlled-rotation-conditional-state}
\end{equation}
である．
\end{prop}
\begin{proof}
式\eqref{eq:controlled-rotation-block}は，互いに直交する部分空間
$\operatorname{span}\{\ket z\}\otimes\mathbb C^2$ごとに
ユニタリ$R_y(2\arcsin g(z))$を作用させるブロック対角演算である．実際，
\begin{align*}
 W_g^\dagger W_g
 &=\sum_{z,z'}\ket z\braket{z|z'}\bra{z'}
   \otimes R_y(2\arcsin g(z))^\dagger R_y(2\arcsin g(z'))\\
 &=\sum_z\ket z\bra z\otimes I=I
\end{align*}
である．入力へ作用させた状態は
\[
 \sum_z\alpha_z\sqrt{1-g(z)^2}\ket z\ket0
 +\sum_z\alpha_zg(z)\ket z\ket1
\]
となる．第2項のノルムの二乗を計算基底の正規直交性で求めると
式\eqref{eq:controlled-rotation-success-probability}を得る．
射影後の第2項を$\sqrt{p_1}$で割れば
式\eqref{eq:controlled-rotation-conditional-state}となる．
\end{proof}

HHLでは，$z$が固有値の近似値$\widetilde\lambda$を表し，
$g(\widetilde\lambda)=C/\widetilde\lambda$と選ぶ．
定数$C$の条件，固有値の符号，有限精度による誤差は第4章で定める．
本節では，角度を計算する可逆回路の構成までは扱わず，
式\eqref{eq:controlled-rotation-block}の演算を必要な精度で実行できるものとする．

